% Generated draft: requires mathematical and source review.
\documentclass[11pt]{article}
\usepackage[margin=25mm]{geometry}
\usepackage{fontspec}
\setmainfont{Latin Modern Roman}
\usepackage{amsmath,amssymb,mathtools,unicode-math}
\setmathfont{Latin Modern Math}
\usepackage{enumitem,needspace,xurl,hyperref}
\setlength{\parindent}{0pt}
\setlength{\parskip}{5pt}
\setlength{\emergencystretch}{3em}
\newcommand{\N}{\mathbb{N}}
\newcommand{\Z}{\mathbb{Z}}
\newcommand{\Q}{\mathbb{Q}}
\newcommand{\R}{\mathbb{R}}
\newcommand{\C}{\mathbb{C}}
\newcommand{\F}{\mathbb{F}}
\newcommand{\A}{\mathbb{A}}
\newcommand{\PP}{\mathbb{P}}
\newcommand{\E}{\mathbb{E}}
\newcommand{\eps}{\varepsilon}
\newcommand{\dd}{\,\mathrm{d}}
\begin{document}
\section*{UnsolvedMath ID: 20001149}
\textbf{Source problem code:} catalogue entry
\subsection{Definitions and mathematical statement}
This record concatenates several workshop sections and has no single complete conjecture. Its opening says ``2-join decomposition is $\mathcal C_{\mathrm{pc}}$-preserving,'' without defining decomposition blocks, followed by the conjecture that no induced-minimal non-perfectly-contractile graph has a star cutset. An even pair is a nonadjacent pair whose induced connecting paths all have even length. A graph is perfectly contractile if every induced subgraph can be reduced to a clique by even-pair contractions; write $\mathcal C_{\mathrm{pc}}$ for this hereditary class. A star cutset is a disconnecting vertex set containing a vertex adjacent to all its other members. The source says PCGC implies ``all three'' conjectures, but the third antecedent and full PCGC statement are not included.

For the aid's working non-path $2$-join setting, partition $V=X_1\cup X_2$ with disjoint nonempty attachment sets $A_i,B_i\subseteq X_i$, cross edges exactly the complete joins $A_1$--$A_2$ and $B_1$--$B_2$, and induced attachment paths in each side, with neither whole side such a path. If $P_i$ is an actual induced $A_i$--$B_i$ path with interior outside $A_i\cup B_i$, the honest block $G[X_1\cup V(P_2)]$ is a proper induced subgraph, as is the opposite block. Thus $G\in\mathcal C_{\mathrm{pc}}$ implies these actual-path blocks belong to the class. This forward hereditary observation is not the undefined full marker-block equivalence. The canonical crossing cycle has length $|E(P_1)|+|E(P_2)|+2$; in an odd-hole-free graph the path parities must agree, so replacing a side needs the correct marker parity.

The aid refutes only a marker-free side-only reading. Take triangles on $\{x_1,x_2,x_3\}$ and $\{y_1,y_2,y_3\}$, matching edges $x_iy_i$, and new vertex $z$ adjacent only to $x_1,y_1$. The sides $\{z,x_1,y_1\}$ and $\{x_2,x_3,y_2,y_3\}$ induce $K_3$ and $C_4$, both perfectly contractile; attachments are $\{x_1\}$ to $\{x_2,x_3\}$ and $\{y_1\}$ to $\{y_2,y_3\}$. The six $x/y$ vertices induce $\overline{C_6}$, which has no even pair and is not a clique, so the whole graph is not perfectly contractile. This does not refute a parity-correct marker-path conjecture.

The remaining extracted topics must be retained as separate historical questions. A hole is an induced cycle of length at least four; a Berge graph has no odd hole or odd antihole of length at least five, and a perfect graph has $\chi(H)=\omega(H)$ for every induced subgraph $H$.
\begin{itemize}
\item F.5 proposes that every even-pair-free Berge graph has, in it or its complement, a bipartite line-graph or $2$-join structure. The grammatical scope of the $2$-join alternative needs confirmation. The source suggests a direct proof might shorten the SPGC proof. It also asks for a sufficient condition on an even pair making $G$ Berge if and only if its contraction is Berge, noting that an arbitrary even-pair contraction can turn a non-Berge graph into a Berge graph.
\item F.6 asks for a polynomial algorithm finding an odd-length $a$--$b$ walk whose entire vertex set induces no triangle; edges may repeat. It reports implications for odd-hole finding and even-pair detection in odd-hole-free graphs, Bienstock's NP-hardness of finding an odd hole through a specified vertex, and the asserted equivalence in a $2$-connected graph between a triangle-free odd closed walk based at $a$ and an odd hole somewhere, not necessarily through $a$.
\item G.1 asks whether $2$-divisible is equivalent to odd-hole-free. A $k$-division partitions vertices so no part contains a maximum-cardinality clique; the hereditary definition excludes edgeless induced subgraphs, as clarified later in G.4. The source states perfect graphs are $2$-divisible and odd holes admit no $2$-division.
\item G.2 asks for a constant clique-color bound for perfect graphs, specifically whether three colors always suffice, and whether the same holds for odd-hole-free graphs. A clique coloring forbids monochromatic inclusionwise maximal cliques of size greater than one, not just maximum cliques. The source records two/three-color results for comparability, cocomparability and further perfect classes, and an ``almost all'' perfect-graph result. It calls the hereditary clique-color property strong $k$-divisibility and asks whether strong $2$-divisibility can be recognized in polynomial time. Its word ``subgraphs'' needs an induced-subgraph convention; the source additionally reports hardness of verifying a proposed clique coloring even in particular perfect classes.
\item G.3 asks for polynomial recognition of odd-hole-free graphs. G.4 defines hereditary $k$-divisibility on induced subgraphs with an edge and records a polynomial even-hole-free recognition algorithm. These are historical source questions and attributions, not an assertion that their present-day status is unchanged.
\end{itemize}
\subsection{Short English statement}
The extraction combines an underdefined perfect-contractility decomposition conjecture with several neighboring perfect-graph problems. The supplied aid proves only forward preservation for actual-path blocks and refutes a marker-free interpretation; all adjacent questions, qualifications and attributions remain separately visible.
\subsection{Sources}
\begin{itemize}
\item Local UnsolvedMath registry, permanent ID $20001149$, source code catalogue entry; canonical \texttt{aim-combinatorics-notes.json}, \texttt{notes[273]}. AIM workshop \emph{The Perfect Graph Conjecture}: \url{https://aimath.org/WWN/perfectgraph/perfectgraph.pdf}. The record is reconstructed and unverified; page headers are extraction artifacts and no PDF inspection is claimed here.
\item Opening decomposition/star-cutset conjectures: Claudia Linhares Sales. F.5 structural proposal: Robin Thomas; even-pair contraction question: Paul Seymour. F.6 walk question: Andr\'as Seb\H{o} and Nicolas Trotignon, with a stated Bienstock hardness result. G.1 question contributed by Bruce Reed, attributed to Hoang--McDiarmid. G.2 contributed by Myriam Preissmann and Andr\'as Seb\H{o}. G.3 recognition question: Chinh Hoang.
\item Actual supplied G.2 references: Duffus et al., source label [92e:06009]; Bacs\'o et al., no identifier; Pr\"omel--Steger, source citation ``Probability and Computation, 1, 1992''; Hoang--McDiarmid, [2002j:05110]. Full titles and exact statements are not supplied. G.4 credits Conforti, Cornu\'ejols, Kapoor and Vu\v{s}kovi\'c with even-hole-free recognition.
\item Attached research attempt 1: invalid statement, recorded run status unknown, rigor label proved, novelty confidence low; registry status open, classification OPEN-TRIAGE, suggested difficulty L2. Its repaired block observations and seven-vertex example are machine generated and not peer reviewed; independently verify before citation.
\end{itemize}
\end{document}
